\bibitem{zurek2009} W. H. Zurek, Quantum Darwinism, \emph{Nat. Phys.} \textbf{5}, 181 (2009).
\bibitem{zurek2022} W. H. Zurek, Quantum theory of the classical: Einselection, envariance, Quantum Darwinism and extantons, \emph{Entropy} \textbf{24}, 1520 (2022).
\bibitem{zurek2025book} W. H. Zurek, \emph{Decoherence and Quantum Darwinism: From Quantum Foundations to Classical Reality} (Cambridge University Press, Cambridge, 2025).
\bibitem{korbicz2014} J. K. Korbicz, P. Horodecki, and R. Horodecki, Objectivity in a noisy photonic environment through quantum state information broadcasting, \emph{Phys. Rev. Lett.} \textbf{112}, 120402 (2014).
\bibitem{horodecki2015} R. Horodecki, J. K. Korbicz, and P. Horodecki, Quantum origins of objectivity, \emph{Phys. Rev. A} \textbf{91}, 032122 (2015).
\bibitem{brandao2015} F. G. S. L. Brand\~ao, M. Piani, and P. Horodecki, Generic emergence of classical features in quantum Darwinism, \emph{Nat. Commun.} \textbf{6}, 7908 (2015).
\bibitem{le2019} T. P. Le and A. Olaya-Castro, Strong Quantum Darwinism and strong independence are equivalent to spectrum broadcast structure, \emph{Phys. Rev. Lett.} \textbf{122}, 010403 (2019).
\bibitem{korbicz2021} J. K. Korbicz, Roads to objectivity: Quantum Darwinism, spectrum broadcast structures, and strong quantum Darwinism--a review, \emph{Quantum} \textbf{5}, 571 (2021).
\bibitem{fu2021} H.-F. Fu, Uniqueness of the observable leaving redundant imprints in the environment in the context of Quantum Darwinism, \emph{Phys. Rev. A} \textbf{103}, 042210 (2021).
\bibitem{touil2022} A. Touil, B. Yan, D. Girolami, S. Deffner, and W. H. Zurek, Eavesdropping on the decohering environment: Quantum Darwinism, amplification, and the origin of objective classical reality, \emph{Phys. Rev. Lett.} \textbf{128}, 010401 (2022).
\bibitem{girolami2022} D. Girolami, A. Touil, B. Yan, S. Deffner, and W. H. Zurek, Redundantly amplified information suppresses quantum correlations in many-body systems, \emph{Phys. Rev. Lett.} \textbf{129}, 010401 (2022).
\bibitem{touil2024branch} A. Touil, F. Anza, S. Deffner, and J. P. Crutchfield, Branching states as the emergent structure of a quantum universe, \emph{Quantum} \textbf{8}, 1494 (2024).
\bibitem{touil2025} A. Touil, B. Yan, and W. H. Zurek, Consensus about classical reality in a quantum universe, arXiv:2503.14791 (2025).
\bibitem{doucet2025} E. Doucet and S. Deffner, Compatibility of quantum measurements and the emergence of classical objectivity, \emph{Phys. Rev. A} \textbf{111}, 042217 (2025).
\bibitem{kiely2026} A. Kiely, D. A. Chisholm, A. Touil, S. Deffner, G. Landi, and S. Campbell, Metrological approach to the emergence of classical objectivity, \emph{Phys. Rev. A} \textbf{113}, 022403 (2026).
\bibitem{torvinen2026} J. Torvinen, E. Keski-Vakkuri, and N. Pranzini, Quantum Darwinism and the quality of Petz recovery, arXiv:2605.06848 (2026).
\bibitem{alfsen2012} E. Alfsen and F. Shultz, Finding decompositions of a class of separable states, arXiv:1202.3673 (2012).
\bibitem{allman2009} E. S. Allman, C. Matias, and J. A. Rhodes, Identifiability of parameters in latent structure models with many observed variables, \emph{Ann. Stat.} \textbf{37}, 3099 (2009).
\bibitem{anandkumar2014} A. Anandkumar, R. Ge, D. Hsu, S. M. Kakade, and M. Telgarsky, Tensor decompositions for learning latent variable models, \emph{J. Mach. Learn. Res.} \textbf{15}, 2773 (2014).
\bibitem{bhaskara2014} A. Bhaskara, M. Charikar, and A. Vijayaraghavan, Uniqueness of tensor decompositions with applications to polynomial identifiability, in \emph{Proc. 26th ACM-SIAM Symposium on Discrete Algorithms} (2014), arXiv:1304.8087.
\bibitem{sidiropoulos2017} N. D. Sidiropoulos, L. De Lathauwer, X. Fu, K. Huang, E. E. Papalexakis, and C. Faloutsos, Tensor decomposition for signal processing and machine learning, \emph{IEEE Trans. Signal Process.} \textbf{65}, 3551 (2017).
\bibitem{helstrom1976} C. W. Helstrom, \emph{Quantum Detection and Estimation Theory} (Academic Press, New York, 1976).
\bibitem{fuchscaves1996} C. A. Fuchs and C. M. Caves, Mathematical techniques for quantum communication theory, \emph{Open Syst. Inf. Dyn.} \textbf{3}, 345 (1995); arXiv:quant-ph/9604001.
\bibitem{fuchsvandegraaf1999} C. A. Fuchs and J. van de Graaf, Cryptographic distinguishability measures for quantum-mechanical states, \emph{IEEE Trans. Inf. Theory} \textbf{45}, 1216 (1999).
\bibitem{huang2020} H.-Y. Huang, R. Kueng, and J. Preskill, Predicting many properties of a quantum system from very few measurements, \emph{Nat. Phys.} \textbf{16}, 1050 (2020).
\bibitem{stephens2021} A. Stephens, J. M. Cutshall, T. McPhee, and M. Beck, Self-consistent state and measurement tomography with fewer measurements, \emph{Phys. Rev. A} \textbf{104}, 012416 (2021).
\bibitem{cattaneo2023} M. Cattaneo, M. A. C. Rossi, K. Korhonen, E.-M. Borrelli, G. Garc\'ia-P\'erez, Z. Zimbor\'as, and D. Cavalcanti, Self-consistent quantum measurement tomography based on semidefinite programming, \emph{Phys. Rev. Research} \textbf{5}, 033154 (2023).
\bibitem{chen2019} M.-C. Chen, H.-S. Zhong, Y. Li, D. Wu, X.-L. Wang, L. Li, N.-L. Liu, C.-Y. Lu, and J.-W. Pan, Emergence of classical objectivity of Quantum Darwinism in a photonic quantum simulator, \emph{Sci. Bull.} \textbf{64}, 580 (2019).
\bibitem{zhu2025} Z. Zhu \emph{et al.}, Observation of quantum Darwinism and the origin of classicality with superconducting circuits, \emph{Sci. Adv.} \textbf{11}, eadx6857 (2025).
\bibitem{galve2016} F. Galve, R. Zambrini, and S. Maniscalco, Non-Markovianity hinders Quantum Darwinism, \emph{Sci. Rep.} \textbf{6}, 19607 (2016).
\bibitem{duruisseau2023} P. Duruisseau, A. Touil, and S. Deffner, Pointer states and Quantum Darwinism with two-body interactions, \emph{Entropy} \textbf{25}, 1573 (2023).

\end{thebibliography}

\clearpage
\appendix

\section{Proof details}
\label{app:proofs}

\subsection{Proof of Theorem~\ref{thm:subspace}}
Vectorize Hermitian operators in a real Hilbert-Schmidt basis and define
\begin{equation}
A_i=[\sqrt{p_1}|\delta_{i,1}\rangle\!\rangle,\ldots,\sqrt{p_r}|\delta_{i,r}\rangle\!\rangle],
\end{equation}
with analogous $A_j$.  Eq.~\eqref{eq:multi-connected} corresponds to $A_iA_j^\top$.  Because $A_i\sqrt p=A_j\sqrt p=0$, where $\sqrt p=(\sqrt{p_1},\ldots,\sqrt{p_r})^\top$, and both matrices have rank $r-1$, each nullspace is exactly $\Span\{\sqrt p\}$.  Therefore $\operatorname{range}(A_j^\top)=(\sqrt p)^\perp$, and the restriction of $A_i$ to that subspace maps onto $\operatorname{range}(A_i)=\cS_i$.  Hence $\operatorname{range}(A_iA_j^\top)=\cS_i$.  The right-support statement is symmetric.

\subsection{Finite-shot concentration for Theorem~\ref{thm:finite-binary}}
Within each physical realization, $(R_i,R_j)$ is a same-event snapshot pair, so $\mathbb E[R_i\otimes R_j]$ retains the desired two-fragment correlation. Let $(R_i',R_j')$ be an independent second realization. Then
\begin{align}
&\mathbb E[(R_i-R_i')\otimes(R_j-R_j')]\notag\\
&\qquad=2\mathbb E[R_i\otimes R_j]-2\mathbb E[R_i]\otimes\mathbb E[R_j],
\end{align}
which proves unbiasedness after the factor $1/2$. Set
\begin{equation}
K_s=\frac12(R_i^{2s-1}-R_i^{2s})\otimes(R_j^{2s-1}-R_j^{2s}).
\end{equation}
The snapshot bounds imply
\begin{equation}
\|K_s\|_2\le \frac12(2B_i)(2B_j)=2B_iB_j=:L.
\end{equation}
For the Hilbert-space empirical mean $Z=N^{-1}\sum_s(K_s-\mathbb EK_s)$, independence gives
\begin{equation}
\mathbb E\|Z\|_2^2
=\frac1{N^2}\sum_s\mathbb E\|K_s-\mathbb EK_s\|_2^2
\le\frac{L^2}{N},
\end{equation}
so $\mathbb E\|Z\|_2\le L/\sqrt N$. Replacing one pair changes $\|Z\|_2$ by at most $2L/N$. McDiarmid therefore yields, with probability at least $1-\beta$,
\begin{equation}
\|Z\|_2\le\frac{L}{\sqrt N}
\left[1+\sqrt{2\ln(1/\beta)}\right],
\end{equation}
which is Eq.~\eqref{eq:finite-omega} because $L=2B_iB_j$. Wedin's rank-one bound gives
\begin{equation}
\sin\theta\le\frac{e_N}{\lambda-e_N}.
\end{equation}
Requiring this to be at most $\varepsilon$ is equivalent to $e_N\le\varepsilon\lambda/(1+\varepsilon)$. Substituting $\lambda=\|\Delta_i\|_2\|\Delta_j\|_2/4$ and converting from $N$ independent pairs to $n=2N$ physical events gives the factor $128$ in Eq.~\eqref{eq:finite-binary-n}. This independently checks both the bounded-difference constant and the pair-to-event conversion.

\subsection{Proof of Theorem~\ref{thm:response-conditioning}}
